%----------------------------------------------------------------------------
\appendix
%----------------------------------------------------------------------------
\chapter*{\fuggelek}\addcontentsline{toc}{chapter}{\fuggelek}
\setcounter{chapter}{\appendixnumber}
%\setcounter{equation}{0} % a fofejezet-szamlalo az angol ABC 6. betuje (F) lesz
\numberwithin{equation}{section}
\numberwithin{figure}{section}
\numberwithin{lstlisting}{section}
%\numberwithin{tabular}{section}

%----------------------------------------------------------------------------
\section{Parameter counts}
\label{app:parameters}
%----------------------------------------------------------------------------

This appendix details the parameter counts reported in \autoref{sec:prefix-adapter} and \autoref{sec:model-size}. A linear layer with a bias, mapping $a$ input features to $b$ output features, has $ab+b$ parameters. The adapter's layer normalisation over $d$ features has $2d$ parameters: one learned scale and one learned shift per feature. A rank-$r$ LoRA update adds $r(a+b)$ parameters.

\autoref{tab:params-trainable} lists the trainable components. In the GNN estimator, the four heads of the GATv2 layer each have two transformations from 1024 to 256 values (\autoref{eq:gnn-head}). In LoRA, the query projection maps 5120 values to 5120, while the value projection maps 5120 values to the $8 \times 128 = 1024$ values of the eight key--value heads.

\begin{table}[!ht]
  \centering
  \small
  \begin{tabular}{@{}llr@{}}
    \toprule
    part & calculation & parameters \\
    \midrule
    \multicolumn{3}{@{}l}{\textit{GNN estimator}} \\
    head transformations $\mathbf{W}^{(k)}_1, \mathbf{W}^{(k)}_2$, with biases & $2 \times 4 \times (1024 \times 256 + 256)$ & 2,099,200 \\
    attention vectors $\mathbf{a}^{(k)}$ & $4 \times 256$ & 1,024 \\
    bias $\mathbf{b}$ of the GATv2 layer & $1024$ & 1,024 \\
    \texttt{[MASK]} vector $\mathbf{m}$ & $1024$ & 1,024 \\
    linear output layer & $1024 \times 1024 + 1024$ & 1,049,600 \\
    total &  & 3,151,872 \\
    \midrule
    \multicolumn{3}{@{}l}{\textit{prefix adapter}} \\
    trunk: linear layer & $1024 \times 1024 + 1024$ & 1,049,600 \\
    trunk: layer normalisation & $2 \times 1024$ & 2,048 \\
    code projection & $1024 \times 15{,}360 + 15{,}360$ & 15,744,000 \\
    two shared key and value decoders & $2 \times (192 \times 32{,}768 + 32{,}768)$ & 12,648,448 \\
    learned base prefix $\mathbf{B}$ & $40 \times 2 \times 8 \times 32 \times 128$ & 2,621,440 \\
    gates $g$ & $40 \times 2$ & 80 \\
    total &  & 32,065,616 \\
    \midrule
    \multicolumn{3}{@{}l}{\textit{LoRA (rank 16)}} \\
    query projections & $40 \times 16 \times (5120 + 5120)$ & 6,553,600 \\
    value projections & $40 \times 16 \times (5120 + 1024)$ & 3,932,160 \\
    total &  & 10,485,760 \\
    \midrule
    all trainable components &  & 45,703,248 \\
    \bottomrule
  \end{tabular}
  \caption{Parameter counts of the trainable components.}
  \label{tab:params-trainable}
\end{table}

The direct mapping considered in \autoref{sec:prefix-adapter} would use 80 linear projections with biases, one from 1024 to $8 \times 32 \times 128 = 32{,}768$ values for the keys and one for the values of each of the 40 layers. This gives $80 \times (1024 \times 32{,}768 + 32{,}768) = 2{,}686{,}976{,}000$ parameters, about 84 times the 32,065,616 parameters of the factorised adapter. Calculating the parameter count from the Qwen3-14B configuration gives 14,768,307,200 parameters~\cite{qwen2025config}. The direct mapping alone would therefore have about 18\% as many parameters as the frozen backbone.

%----------------------------------------------------------------------------
\section{Generation prompt}
\label{app:generation-prompt}
%----------------------------------------------------------------------------

Listing~\ref{lst:generation-prompt} gives the complete fixed prompt of the reader (\autoref{sec:reader}). It consists of a system message and a user message, formatted with Qwen3's chat template with the reasoning mode disabled; it is identical for every paper.

\begin{lstlisting}[caption={Fixed prompt of the reader.},label={lst:generation-prompt}]
System:
You are a scientific writing assistant. The context you have been given encodes the reference list of one paper, whose own abstract is missing. That paper is not among its references: it is a separate, more recent work that cites all of them. Reconstruct its abstract as continuous prose covering, in order: the problem it addresses; what the work it cites has already established about that problem; what it does about what remains open, described by the mechanism it uses the kind of model, data or analysis rather than by a name; and what it evaluates. Where the context leaves the paper's specific contribution undetermined, state what it does determine about its problem and approach, and stop there; a correct general statement is worth more than a specific guess. Write in the voice of the paper's own authors: never mention the cited works as a collection, how many there are, or the authors, and do not restate any single one of them. Write plain prose of 150-180 words, with no markup, mathematical notation or headings. Output ONLY the abstract text.

User:
Write the abstract.
\end{lstlisting}

%----------------------------------------------------------------------------
\section{Judge prompts}
\label{app:judge-prompt}
%----------------------------------------------------------------------------

This section gives the system prompts of the two judge calls (\autoref{sec:judge-evaluation}). In the content call (Listing~\ref{lst:judge-content}), the user message contains the original abstract after the line ``ORIGINAL ABSTRACT:'' and the generated abstract after the line ``GENERATED ABSTRACT:''. In the quality call (Listing~\ref{lst:judge-quality}), the user message contains only the generated abstract, after the line ``ABSTRACT:''.

\begin{lstlisting}[caption={System prompt of the content call.},label={lst:judge-content}]
You are a rigorous scientific peer reviewer. Given an original paper abstract and a generated abstract for the same paper, compare them. Return valid JSON with exactly this structure:
{
  "subfield":            <int 0-3>,
  "problem":             <int 0-3>,
  "task":                <int 0-3>
}

subfield, problem and task use one scale:
  0 = different, 1 = related but clearly different, 2 = the same with differences in detail, 3 = the same.
- subfield: the specific research niche, narrower than the field (e.g. "graph contrastive learning", not "machine learning").
- problem: the underlying problem, gap or motivation the paper addresses.
- task: the concrete task or objective the paper takes on (e.g. detection vs segmentation, translation vs summarisation).
Do not let the methods, models, datasets or numbers named in either abstract influence these three scores.
\end{lstlisting}

\begin{lstlisting}[caption={System prompt of the quality call.},label={lst:judge-quality}]
You are a rigorous scientific peer reviewer. Rate the writing of the abstract below, without judging whether its content is correct. Return valid JSON with exactly this structure:
{
  "coherence": <int 0-3>,
  "language":  <int 0-3>
}

- coherence: does the abstract form one logical argument?
  0 = no recognisable argument; 1 = a list of loosely connected statements; 2 = a clear argument with gaps or jumps; 3 = a logical progression from the problem to the approach to the results.
- language: does the text follow the norms of scientific writing?
  0 = not scientific prose; 1 = frequent generic or conversational phrasing; 2 = scientific register with some vague or imprecise terms; 3 = precise domain terminology throughout, with no generic filler.
\end{lstlisting}

%----------------------------------------------------------------------------
\section{Training resources}
\label{app:resources}
%----------------------------------------------------------------------------

All training ran on the Komondor supercomputer, with one NVIDIA A100 GPU with 40\,GiB of memory per job. The weights of Qwen3-14B are loaded in bfloat16, while the LoRA parameters and the loss are kept in float32, and gradient checkpointing is enabled during training. The implementation uses transformers 4.51.3, pinned exactly because inserting the prefix into the key--value cache depends on internals of the library that change between versions. \autoref{tab:resources} lists the training time of each phase.

\begin{table}[!ht]
  \centering
  \begin{tabular}{@{}lr@{}}
    \toprule
    phase & wall-clock time \\
    \midrule
    GNN estimator, 8 variants & 31\,min in total \\
    reader & 11.5\,h \\
    joint fine-tuning & 5.7\,h \\
    \bottomrule
  \end{tabular}
  \caption{Training time of each phase of the final model on one A100 GPU. The times include model loading and the validation during training.}
  \label{tab:resources}
\end{table}

%----------------------------------------------------------------------------
\section{Example generations}
\label{app:examples}
%----------------------------------------------------------------------------

This section shows the original and the generated abstract of four evaluation papers, generated by the final model (\autoref{sec:results-final}). I chose the papers to illustrate one case of stronger content agreement and three distinct limitations; they are not a random sample and do not show how often each case occurs. The abstracts are reproduced as recorded in the dataset and as generated, without editing.

\subsection{Stronger content agreement}
\textit{Dense dilated convolutions merging network for land cover classification}~\cite{liu2020ddcmnet}. Judge scores: subfield 3, problem 3, task 3, coherence 3, language 2.\nopagebreak[4]

\textbf{Original abstract.} Land cover classification of remote sensing images is a challenging task due to limited amounts of annotated data, highly imbalanced classes, frequent incorrect pixel-level annotations, and an inherent complexity in the semantic segmentation task. In this article, we propose a novel architecture called the dense dilated convolutions' merging network (DDCM-Net) to address this task. The proposed DDCM-Net consists of dense dilated image convolutions merged with varying dilation rates. This effectively utilizes rich combinations of dilated convolutions that enlarge the network's receptive fields with fewer parameters and features compared with the state-of-the-art approaches in the remote sensing domain. Importantly, DDCM-Net obtains fused local- and global-context information, in effect incorporating surrounding discriminative capability for multiscale and complex-shaped objects with similar color and textures in very high-resolution aerial imagery. We demonstrate the effectiveness, robustness, and flexibility of the proposed DDCM-Net on the publicly available ISPRS Potsdam and Vaihingen data sets, as well as the DeepGlobe land cover data set. Our single model, trained on three-band Potsdam and Vaihingen data sets, achieves better accuracy in terms of both mean intersection over union (mIoU) and F1-score compared with other published models trained with more than three-band data. We further validate our model on the DeepGlobe data set, achieving state-of-the-art result 56.2\% mIoU with much fewer parameters and at a lower computational cost compared with related recent work. Code available at \url{https://github.com/samleoqh/DDCM-Semantic-Segmentation-PyTorch}

\textbf{Generated abstract.} In this paper, we propose a novel deep learning framework for semantic segmentation of remote sensing images. The proposed framework is based on a multi-scale convolutional neural network (CNN) with a novel multi-scale feature fusion module. The multi-scale feature fusion module is designed to combine features from different scales and to preserve the spatial information of the input image. The proposed framework is trained on a large-scale remote sensing image dataset, which is collected from the Google Earth. The proposed framework is evaluated on the ISPRS Vaihingen dataset and the results show that the proposed framework outperforms the state-of-the-art methods.

This example agrees with the original in research area and task, but the high content scores do not establish the accuracy of its architectural and experimental details.

\subsection{Related area, different task}
\textit{Learn to segment retinal lesions and beyond}~\cite{wei2021retinal}. Judge scores: subfield 1, problem 1, task 1, coherence 3, language 3.\nopagebreak[4]

\textbf{Original abstract.} Towards automated retinal screening, this paper makes an endeavor to simultaneously achieve pixel-level retinal lesion segmentation and image-level disease classification. Such a multi-task approach is crucial for accurate and clinically interpretable disease diagnosis. Prior art is insufficient due to three challenges, that is, lesions lacking objective boundaries, clinical importance of lesions irrelevant to their size, and the lack of one-to-one correspondence between lesion and disease classes. This paper attacks the three challenges in the context of diabetic retinopathy (DR) grading. We propose L-Net, a new variant of fully convolutional networks, with its expansive path re-designed to tackle the first challenge. A dual loss that leverages both semantic segmentation and image classification losses is devised to resolve the second challenge. We propose Side-Attention Net (SiAN) as our multi-task framework. Harnessing L-Net as a side-attention branch, SiAN simultaneously improves DR grading and interprets the decision with lesion maps. A set of 12K fundus images is manually segmented by 45 ophthalmologists for 8 DR-related lesions, resulting in 290K manual segments in total. Extensive experiments on this large-scale dataset show that our proposed approach surpasses the prior art for multiple tasks including lesion segmentation, lesion classification and DR grading.

\textbf{Generated abstract.} We propose a novel deep learning architecture for the task of retinal vessel segmentation. The proposed architecture is based on a multi-scale convolutional neural network (CNN) that is trained to predict the probability of a pixel being a vessel. The multi-scale CNN is trained on a dataset of 100 retinal images and tested on a separate dataset of 20 retinal images. The proposed architecture is compared with a state-of-the-art CNN architecture and a multi-scale CNN architecture. The proposed architecture achieves a mean intersection over union (IoU) of 0.72, which is 1.5\% higher than the state-of-the-art CNN and 0.5\% higher than a multi-scale CNN.

\subsection{Analysis recast as a proposed method}
\textit{Analysis of gender inequality in face recognition accuracy}~\cite{albiero2020gender}. Judge scores: subfield 3, problem 1, task 3, coherence 3, language 2.\nopagebreak[4]

\textbf{Original abstract.} We present a comprehensive analysis of how and why face recognition accuracy differs between men and women. We show that accuracy is lower for women due to the combination of (1) the impostor distribution for women having a skew toward higher similarity scores, and (2) the genuine distribution for women having a skew toward lower similarity scores. We show that this phenomenon of the impostor and genuine distributions for women shifting closer towards each other is general across datasets of African-American, Caucasian, and Asian faces. We show that the distribution of facial expressions may differ between male/female, but that the accuracy difference persists for image subsets rated confidently as neutral expression. The accuracy difference also persists for image subsets rated as close to zero pitch angle. Even when removing images with forehead partially occluded by hair/hat, the same impostor/genuine accuracy difference persists. We show that the female genuine distribution improves when only female images without facial cosmetics are used, but that the female impostor distribution also degrades at the same time. Lastly, we show that the accuracy difference persists even if a state-of-the-art deep learning method is trained from scratch using training data explicitly balanced between male and female images and subjects.

\textbf{Generated abstract.} Face recognition is a challenging task in the wild due to the large variations in illumination, pose, and expression. In this paper, we propose a novel face recognition method that is robust to these variations. Our method is based on a deep convolutional neural network (CNN) that is trained to predict the identity of a face image. We propose a new loss function that is based on the cosine similarity between the predicted and ground-truth identity vectors. We also propose a new training procedure that is based on a combination of hard and easy samples. We show that our method achieves state-of-the-art performance on the LFW, IJB-A, and IJB-B face recognition benchmarks.

Despite the high task score, the generated abstract changes the paper's contribution from analysing accuracy differences to proposing a recognition method.

\subsection{Description of a recorded reference}
\textit{Exemplar normalization for learning deep representation}~\cite{zhang2020exemplar}. Judge scores: subfield 0, problem 0, task 0, coherence 3, language 3.\nopagebreak[4]

\textbf{Original abstract.} Normalization techniques are important in different advanced neural networks and different tasks. This work investigates a novel dynamic learning-to-normalize (L2N) problem by proposing Exemplar Normalization (EN), which is able to learn different normalization methods for different convolutional layers and image samples of a deep network. EN significantly improves flexibility of the recently proposed switchable normalization (SN), which solves a static L2N problem by linearly combining several normalizers in each normalization layer (the combination is the same for all samples). Instead of directly employing a multi-layer perceptron (MLP) to learn data-dependent parameters as conditional batch normalization (cBN) did, the internal architecture of EN is carefully designed to stabilize its optimization, leading to many appealing benefits. (1) EN enables different convolutional layers, image samples, categories, benchmarks, and tasks to use different normalization methods, shedding light on analyzing them in a holistic view. (2) EN is effective for various network architectures and tasks. (3) It could replace any normalization layers in a deep network and still produce stable model training. Extensive experiments demonstrate the effectiveness of EN in a wide spectrum of tasks including image recognition, noisy label learning, and semantic segmentation. For example, by replacing BN in the ordinary ResNet50, improvement produced by EN is 300\% more than that of SN on both ImageNet and the noisy WebVision dataset.

\textbf{Generated abstract.} We propose a new convolutional neural network (CNN) architecture, called Residual Network (ResNet), that consists of very deep feed-forward networks with the building blocks of ``identity shortcuts'' and ``bottleneck layers''. We apply it to the 1,000-class ImageNet model, where it improves relative error rate by 37.7\% over its predecessor. We then compare it with the state-of-the-art solution on CIFAR-10/100, which shows 6.49\% and 17.0\% improvement.

The paper's recorded references include \textit{Deep Residual Learning for Image Recognition}~\cite{he2016resnet}, the paper that introduced ResNet. The generated abstract describes ResNet rather than exemplar normalization. Its specific numerical claims are not stated in either recorded abstract.

\input{include/\AIdeclaration}
